\section{Flow Matching 目标函数}

\subsection{朴素方法：全局最大似然训练}

一种直觉方法是直接最大化数据的对数似然：

$$\max_\theta \; \mathbb{E}_{x_1 \sim p_{\text{data}}} \left[ \log p_1(x_1) \right]$$

具体来说：
\begin{enumerate}
    \item 用神经网络 $v_\theta(t, x)$ 参数化向量场
    \item 对每个数据点 $x_1$，从 $x_1$ 出发反向运行 ODE 到 $x_0$
    \item 计算 $\log p_0(x_0) - \int_0^1 \nabla \cdot v_\theta(t, \phi_t) \, dt$
    \item 最大化这个值
\end{enumerate}

\begin{warningbox}{朴素方法的致命缺陷}
虽然这个方法在理论上是正确的，但在实践中存在严重问题：
\begin{enumerate}
    \item \textbf{计算代价高}：每次计算损失函数都需要求解 ODE，这本身就是一个昂贵的操作
    \item \textbf{数值不稳定}：ODE 求解器在长时间积分时容易产生数值误差累积
    \item \textbf{反向传播困难}：梯度需要通过整个 ODE 求解过程反向传播（需要 adjoint method），进一步增加计算成本和不稳定性
\end{enumerate}
这就是早期 Continuous Normalizing Flows（如 Neural ODE，Chen et al., 2018）面临的核心困难。
\end{warningbox}

\subsection{Flow Matching 目标：回归向量场}

Flow Matching 的核心突破在于：\textbf{绕过 ODE 求解}，直接用简单的回归损失来训练向量场。

Flow Matching 的损失函数定义为：

$$\mathcal{L}_{\text{FM}}(\theta) = \mathbb{E}_{t \sim \mathcal{U}[0,1], \, x_t \sim p_t} \left\| v_\theta(t, x_t) - v_t(x_t) \right\|^2$$

其中各符号的含义为：
\begin{itemize}
    \item $v_\theta(t, x_t)$：神经网络预测的向量场
    \item $v_t(x_t)$：目标向量场（ground truth）
    \item $t \sim \mathcal{U}[0,1]$：均匀采样时间步
    \item $x_t \sim p_t$：从时刻 $t$ 的边际分布中采样
\end{itemize}

\begin{importantbox}{Flow Matching 的核心优势}
这个损失函数是一个简单的\textbf{均方误差回归}，不需要求解任何 ODE！但问题在于：我们并不知道边际向量场 $v_t(x_t)$ 和边际分布 $p_t(x_t)$ 的解析形式——因为数据分布 $p_{\text{data}}$ 本身就是未知的。

这就引出了 Flow Matching 的关键创新：\textbf{Conditional Flow Matching}。
\end{importantbox}

\subsection{本章小结}

直接最大化似然需要昂贵的 ODE 求解，而 Flow Matching 提出了简洁的回归目标函数。但边际量的不可解析性要求我们进一步引入条件版本的公式。


